\section{Related Work}

\subsection{State-Space Models}

\subsubsection{Classical SSMs}

\textbf{S4} (Gu et al., 2022) introduced structured state-space models with 
diagonal plus low-rank A matrices, achieving $\mathcal{O}(L)$ complexity. 
\textbf{S5} (Smith et al., 2023) parallelized the SSM scan across time.

\subsubsection{Selective SSMs (Mamba)}

\textbf{Mamba} (Gu \& Dao, 2023) introduced \textbf{input-dependent} state transitions:

\[
h_t = \bar{A}(x_t) h_{t-1} + \bar{B}(x_t) x_t
\tag{2.1}
\]

where $\bar{A}, \bar{B}$ depend on the input $x_t$. This selective mechanism 
makes the SSM behave like a learned, content-aware recurrence.

\textbf{Our distinction.} Our implementation \textbf{does not} perform a true selective 
scan. We use learned exponential decay as a \textbf{proxy} for the selectivity, 
implemented as a pairwise softmax over positions.

\subsection{Efficient Attention}

\textbf{Linformer} (Wang et al., 2020) projects K, V to a fixed rank. 
\textbf{Performer} (Choromanski et al., 2021) uses random Fourier features. 
\textbf{Longformer} (Beltagy et al., 2020) uses sparse local windows.

\textbf{Our distinction.} We exploit SSM-inspired decay rather than a fixed 
sequence projection or kernel approximation.

\subsection{Attention-Replacement Modules}

\textbf{Hyena} (Poli et al., 2023) uses long convolutions. \textbf{RWKV} (Peng et al., 
2023) uses a linear recurrence with time-decay. \textbf{RetNet} (Sun et al., 2023) 
uses a retention mechanism with decay.

\textbf{Our distinction.} We operate on pretrained attention weights, reusing 
$W_Q, W_K, W_V, W_O$ in the SSM structure.

\subsection{Position of Our Work}

\begin{table}[H]
\centering
\small
\begin{tabular}{lllll}
\toprule
\textbf{Method} & \textbf{Complexity} & \textbf{Reuses pretrained?} & \textbf{Causal} & \textbf{Bidirectional} \\
\midrule
Mamba & $\mathcal{O}(L)$ & No & Yes & Yes (Bi-Mamba) \\
Hyena & $\mathcal{O}(L \log L)$ & No & Yes & Yes \\
RWKV & $\mathcal{O}(L)$ & No & Yes & No \\
\textbf{Ours} & \textbf{$\mathcal{O}(L^2 D)$} & \textbf{Yes} & \textbf{Yes} & \textbf{Yes} \\
\bottomrule
\end{tabular}
\end{table}

\textbf{Key distinction.} We reuse pretrained attention weights in an SSM-inspired 
formulation, trading complexity for compatibility with pretrained transformers.



\subsection{Comparison Table}

\begin{lstlisting}

---
\end{lstlisting}